\documentclass{article}

\usepackage{amsmath}
\usepackage{amsfonts}
\usepackage{amsthm}
% \usepackage{setspace}
\setlength{\parskip}{1.6em}
\setlength{\parindent}{0pt}
\newtheorem{theorem}{Theorem}
\newtheorem{lemma}{Lemma}
\newtheorem{corollary}{Corollary}

\newcommand{\0}{\mathbf{0}}
\renewcommand{\vec}{\mathbf}
\newcommand{\vv}{\vec{v}}
% \newcommand{\Bdiag}{\text{Bdiag}}
% \newcommand{\pen}[2]{\lambda#1+#2}
% \newcommand{\penEA}{\pen{E}{A}}
\newcommand{\C}{\mathbb{C}}
\newcommand{\N}{\mathbb{N}}
\newcommand{\vecspace}[1]{\mathcal{#1}}
\DeclareMathOperator{\extker}{extker}
\DeclareMathOperator{\id}{id}
\DeclareMathOperator{\img}{Im}
\DeclareMathOperator{\spec}{spec}
% \newcommand{\inv}{^{-1}}

\begin{document}

\section{Notation}
List of notations:

Kernel of $T$ in $V$: $\ker_\vecspace{V}\,T$

Extended kernel of $T$ in $\vecspace{V}$: $\extker T$ is defined as
\begin{equation*}
    \bigcup_{i=0}^{\infty}\ker T^i
\end{equation*}

Identity function for some domain $D$: $\id_D$ or sometimes just $\id$ when $D$ can be implied

Image of function $f$: $\img f$

Spectrum of matrix/linear transformation $T$: $\spec T$ or $\spec_\vecspace{V} T$ for the spectrum in a certain vector space $\vecspace{V}$

For subspaces $\vecspace{S}, \vecspace{S}'$ of $\vecspace{V}$, $\vecspace{S}$ and $\vecspace{S}'$ are disjoint iff $\vecspace{S} \cup \vecspace{S}' = \{\0\}$

For subspaces $\vecspace{S}, \vecspace{S}'$ of $\vecspace{V}$, if $\vecspace{S}$ and $\vecspace{S}'$ are disjoint, we define
\begin{equation*}
    \vecspace{S} \oplus \vecspace{S}' = \{\vec{v}+\vec{u}|\vec{v}\in\vecspace{S},\vec{u}\in\vecspace{S}\}
\end{equation*}
, which is also a subspace of $\vecspace{V}$. \\
Also, the notation $\vecspace{S}\oplus\vecspace{S}'$ imply the two subspaces are disjoint.

A set of finite dimensional subspace $\{\vecspace{S_0}, \vecspace{S_1}, \dots, \vecspace{S_{n-1}}\}$ are independent if the union of their basis are linear independent for any choices of basis.

\section{relevent proof}


Let $\vecspace{V}$ be a $n$ dimensional vector space in $\C^n$, $T:\vecspace{V} \to \vecspace{V}$ be a linear map.
\begin{lemma}
    \label{lemma:extker_is_subspace}
    $\vecspace{S} = \extker_\vecspace{V} T$ is a subspace of $\vecspace{V}$.
\end{lemma}
\begin{proof}
    Note that $\vecspace{S} \subseteq \vecspace{V}$.
    Let $\vec{v}, \vec{u} \in \vecspace{S}$, $k \in \C$. Then, $\exists n,m \in \N$ s.t. $T^{n}\vec{v}=\0$, $T^{m}\vec{u}=\0$.
    % max will suffices, but sum is easier
    Then
    \begin{align*}
        T^{n+m}(\vec{v}+k\vec{u})
        &=T^{n+m}\vec{v}+kT^{n+m}\vec{u}\\
        &=T^m(T^n\vec{v})+kT^n(T^m\vec{u})\\
        &=T^m\0+kT^n\0\\
        &=\0\\
    \end{align*}
    Hence, $\vec{v}+k\vec{u} \in \vecspace{S}$.
\end{proof}

\begin{lemma}
    \label{lemma:ext_finite_power}
    For all $n$ dimensional vector space $\vecspace{V}$ and linear map $T: \vecspace{V} \to \vecspace{V}$ where $n\in\N$, there exists $m\in\N$ s.t. $\extker T = \ker T^m$
\end{lemma}
\begin{proof}
    Note that $\ker T^i\in\ker T^{i+1}$ for all $i\in\N$, so $\dim\ker T^i\le\dim\ker T^{i+1}$.

    Suppose $\ker T^i = \ker T^{i+1}$ for some $i\in\N$, then $\forall \vec{u}\in\vecspace{V}$, $T^i\vec{u}=\0$ iff $T^{i+1}\vec{u}=\0$.
    Hence,
    \begin{align*}
        \ker T^{i+2}
        &= \{\vec{u}\in\vecspace{V}|T^{i+2}\vec{u}=\0\}\\
        &= \{\vec{u}\in\vecspace{V}|TT^{i+1}\vec{u}=\0\}\\
        &= \{\vec{u}\in\vecspace{V}|TT^iu\vec{=}\0\}\\
        &= \{\vec{u}\in\vecspace{V}|TT^i\vec{u}=\0\}\\
        &= \{\vec{u}\in\vecspace{V}|T^{i+1}\vec{u}=\0\}\\
        &= \ker T^{i+1}\\
    \end{align*}
    blah blah blah, you see the point
\end{proof}

\begin{lemma}
    \label{lemma:independent_from_vectors}
    Subspaces $\vecspace{S}_0, \vecspace{S}_1, \dots, \vecspace{S}_{n-1}$ are independent iff for all $\vec{v}_0\in\vecspace{S}_0, \vec{v}_1\in\vecspace{S}_1, \dots, \vec{v}_{n-1}\in\vecspace{S}_{n-1}$, if $\vec{v}_i\neq\0$ for all $i\in[n]$, then $\vec{v}_0, \vec{v}_1, \dots, \vec{v}_{n-1}$ are linearly independent.
\end{lemma}
\begin{proof}
    Let $\beta_0, \beta_1, \dots, \beta_{n-1}$ be any basis of the subspaces.

    To show the forward direction, suppose $\vec{v}_0x_0+\vec{v}_1x_1+\dots+\vec{v}_{n-1}x_{n-1}=\0$, then $\vec{v}_0x_0+\vec{v}_1x_1+\dots+\vec{v}_{n-1}x_{n-1}=\0$.
    And $\beta_0\lfloor\vec{v}_0\rfloor_{\beta_0}x_0 + \dots + \beta_{n-1}\lfloor\vec{v}_{n-1}\rfloor_{\beta_{n-1}}x_{n-1}=\0$. Hence, $x_0,\dots,x_{n-1}=0$.

    For the backward direction, suppose for basis $\gamma_0, \gamma_1, \dots, \gamma_{n-1}$, where $\gamma_i=\{\vec{b}_{i,0}, \dots, \vec{b}_{i_{d_i-1}}\}$, there exists a $x_{0,0},\dots,x_{0,d_0-1},x_{1,0},\dots,x_{n-1,d_{n-1}-1}$ such that $\vec{b}_{0,0}x_{0,0},\dots,\vec{b}_{0,d_0-1}x_{0,d_0-1},\vec{b}_{1,0}x_{1,0},\dots,\vec{b}_{n-1,d_{n-1}-1}x_{n-1,d_{n-1}-1}=\0$.
    If $x_{i,j}=0$ for all $i,j$, then we are done.
    Otherwise, WLOG, $x_{0,0}\neq0$.
    Let $\vec{v}_i=\vec{b}_{i,0}x_{i,0}+\dots+\vec{b}_{i,d_i-1}x_{i,d_i-1}$.
    $\vec{v}_0\neq\0$ and $\vec{v}_0+\vec{v}_1+\dots+\vec{v}_{n-1}=\0$. For any $i$ s.t. $\vec{v}_i=\0$, we can replace it with $0\vec{v}'_i$ where $vec{v}'_i\in\vecspace{S}_i$. This lead to a contradiction.
\end{proof}

% \begin{lemma}
%     \label{lemma:independent_from_one_basis}
%     Subspaces $\vecspace{S}_0, \vecspace{S}_1, \dots, \vecspace{S}_{n-1}$ are independent iff the union of their basis are linear independent for some choices of basis.
% \end{lemma}
% \begin{proof}
%     The forward direction follows directly from definition.
%     To show the backward direction is true, let $\beta_0, \beta_1, \dots, \beta_{n-1}$ be linearly independent basis of the subspaces.
%     Suppose for basis $\gamma_0, \gamma_1, \dots, \gamma_{n-1}$, there exists
% \end{proof}

\section{Main}

Let the eigenvalues of $T$ be $\lambda_0, \lambda_1, \dots, \lambda_{k-1}$. Let $\vecspace{S}_i=\extker T-\lambda_i\id$.

Note that any polynomials of $T$ commute.

% \begin{lemma}
%     \label{lemma:closed_under_T_sub_eigen}
%     $\forall i\in[k]$, $\vecspace{S}_i$ is closed under $T-\lambda_i\id$.
% \end{lemma}
% \begin{proof}
%     $\forall \vec{u} \in \vecspace{S}_i$, $\exists m \in \N$, s.t. $(T-\lambda_i\id)^m\vec{u}=\0$. Hence, $(T-\lambda_i\id)^m(T-\lambda_i\id)\vec{u}=(T-\lambda_i\id)(T-\lambda_i\id)^m\vec{u}=\0$
% \end{proof}

\begin{lemma}
    \label{lemma:closed_under_T}
    $\forall i\in[k]$, $\vecspace{S}_i$ is closed under $T$.
\end{lemma}
\begin{proof}
    $T\vec{u}=(T-\lambda_i\id)\vec{u} + \lambda_i\id\vec{u}=\lambda_i\vec{u}\in\vecspace{S}_i$
\end{proof}

\begin{corollary}
    \label{corollary:closed_under_p_T}
    $\vecspace{S}_i$ is closed under any polynomials of $T$
\end{corollary}
% Note that due to being a subspace, Lemma \ref{lemma:closed_under_T} imply $\vecspace{S}_i$ is closed under any polynomials of $T$.

However, we can go one step further.
$T-\lambda_i\id$ is a well-defined function on $\vecspace{S}_i$ thanks to \ref{corollary:closed_under_p_T}.

\begin{lemma}
    \label{lemma:bijection}
    $\forall i\in[k], x\in\C$, if $x\neq\lambda_i$, $T-x\id$ is a bijection on $\vecspace{S}_i$
\end{lemma}
\begin{proof}
    \begin{align*}
        \ker_{\vecspace{S}_i}(T-x\id)
        &=\{\vec{v}\in\vecspace{S}_i|(T-x\id)\vec{v}=\0\}\\
        &=\{\vec{v}\in\vecspace{V}|T\vec{v}=x\vec{v} \land (T-\lambda_i\id)^{p_i}\vec{v}=\0\}\\
    \end{align*}
    Note that since $T\vec{v}=x\vec{v}$,
    \begin{equation*}
        (T-\lambda_i\id)\vec{v}=T\vec{v}-\lambda_i\vec{v}=(x-\lambda_i)\vec{v}
    \end{equation*}
    where $x-\lambda_i\neq0$ from the assumption. \\
    Hence,
    \begin{align*}
        \ker_{\vecspace{S}_i}(T-x\id)
        &=\{\vec{v}\in\vecspace{V}|T\vec{v}=x\vec{v} \land (x-\lambda_i)^{p_i}\vec{v}=\0\}\\
        &\subseteq\{\vec{v}\in\vecspace{V}|(x-\lambda_i)^{p_i}\vec{v}=\0\}\\
        &=\{\vec{v}\in\vecspace{V}|\vec{v}=\0\}\\
        &=\{\0\}\\
    \end{align*}
    Since $\vecspace{S}_i$ has finite dimension, $T-x\id$ is a bijection on $\vecspace{S}_i$.
\end{proof}

% \begin{lemma}
%     \label{lemma:disjoint}
%     $\forall i\neq j \in [k]$, $\vecspace{S}_i\cup \vecspace{S}_j = \{\0\}$.
% \end{lemma}
% \begin{proof}
%     Let $\vec{u}\in\vecspace{S}_i\cup\vecspace{S}_j$.
%     Suppose $\vec{u}\neq\0$.
%     Let $m_1 = \min\{m\in\N, (T-\lambda_i\id)^m\vec{u}=\0\}$.
%     Note that $m_1\neq0$.
%     Let $\vec{v}=(T-\lambda_i)^{m_1-1}\vec{u}$.
%     Note that $\vec{v}\neq\0$, $\vec{v}\in\vecspace{S}_j$ and $(T-\lambda_i\id)\vec{v}=\0$.
%     Hence, there exists $m_2\in\N$, s.t. $(T-\lambda_2\id)^{m_2}\vec{v}=\0$.
%     However, $(T-\lambda_2\id)\vec{v}=(T-\lambda_1\id+(\lambda_1-\lambda_2)\id)\vec{v}=(\lambda_1-\lambda_2)\vec{v}$.
%     Hence, $(T-\lambda_2\id)^{m_2}\vec{v}=(\lambda_1-\lambda_2)^{m_2}\vec{v}\neq\0$.
%     By contradiction, $\vec{u}=\0$
% \end{proof}

% Let $p_i$ denote $\min\{p\in\N|\extker(T-\lambda_i\id) = \ker(T-\lambda_i\id)^p\}$, whose existance is shown in Lemma \ref{lemma:ext_finite_power}.

% \begin{theorem}
%     \label{theorem:independent}
%     The subspaces $\vecspace{S}_0, \vecspace{S}_1, \dots, \vecspace{S}_{k-1}$ are disjoint.
% \end{theorem}
% \begin{proof}
%     Suppose for $\0\neq\vec{v}_0\in\vecspace{S}_0, \dots, \0\neq\vec{v}_{k-1}\vecspace{S}_{k-1}$, there exists $x_0,\dots,x_{k-1}$ s.t. $x_0\vec{v_0}+x_1\vec{v}_1+\dots+x_{k-1}\vec{v}_{k-1}$ and $x_i\neq0$ for some $i$.

%     WLOG, let $x_0\neq0$.
%     \begin{align*}
%         \0&=(T-\lambda_1\id)^{p_1}(T-\lambda_2\id)^{p_0}\dots(T-\lambda_{k-1}\id)^{p_{k-1}}\left(\sum_{i=0}^{k-1}x_i\vec{v}_i\right)\\
%         &=(T-\lambda_1\id)^{p_1}\dots(T-\lambda_{k-1}\id)^{p_{k-1}}(x_0\vec{v}_0)+\sum_{i=1}^{k-1}(T-\lambda_1\id)^{p_1}\dots(T-\lambda_{k-1}\id)^{p_{k-1}}(x_i\vec{v}_i)\\
%     \end{align*}
%     Since polynomials of $T$ commute, we can rearrange $(T-\lambda_1\id)^{p_1}\dots(T-\lambda_{k-1}\id)^{p_{k-1}}(x_i\vec{v}_i)$, moving the term $(T-\lambda_i\id)^{p_{i}}$ last.
%     Since $(T-\lambda_i\id)^{p_{i}}(x_i\vec{v_i})=\0$, the above equation imply $(T-\lambda_1\id)^{p_1}\dots(T-\lambda_{k-1}\id)^{p_{k-1}}(x_0\vec{v}_0)=\0$.

%     However, for any $\0\neq\vec{v}\in\vecspace{S}_0$ and $i\neq0$, we have $(T-\lambda_{i}\id)^{p_i}\vec{v}\in\vecspace{S}_0$.
%     Also, since $\vec{v}\neq\0$, from Lemma \ref{lemma:disjoint}, we know $\vec{v}\notin\vecspace{S}_i$.
%     Therefore $(T-\lambda_{i}\id)^{p_i}\vec{v}\neq\0$.
%     Hence, $(T-\lambda_1\id)^{p_1}\dots(T-\lambda_{k-1}\id)^{p_{k-1}}(x_0\vec{v}_0)\neq\0$, leading to contradiction.

%     Therefore, $\vecspace{S}_0, \dots, \vecspace{S}_{k-1}$ are independent from Lemma \ref{lemma:independent_from_vectors}.
% \end{proof}

\begin{lemma}
    \label{lemma:img_sum_ker}
    $\ker(T-\lambda_i\id)^{p_i}\oplus\img(T-\lambda_i\id)^{p_i}=\vecspace{V}$
\end{lemma}
\begin{proof}
    First, we will show that the two vector spaces are disjoint.

    Suppose $\vec{v} \in \img(T-\lambda_i\id)^{p_i}$, then there exist $\vec{u}\in\vecspace{V}$ s.t. $\vec{v}=(T-\lambda_i\id)^{p_i}\vec{u}$.
    Since $\vec{v}\in\ker(T-\lambda_i\id)^{p_i}$,
    \begin{align*}
        (T-\lambda_i\id)^{2p_i}\vec{u}
        &=(T-\lambda_i\id)^{p_i}\left((T-\lambda_i\id)^{p_i}\vec{u}\right)\\
        &=(T-\lambda_i\id)^{p_i}\vec{v}\\
        &=\0
    \end{align*}
    Hence, $\vec{u}\in\extker(T-\lambda_i\id)=\ker(T-\lambda_i\id)^{p_i}$.
    \begin{equation*}
        \vec{v} = (T-\lambda_i\id)^{p_i}\vec{u} = \0
    \end{equation*}
    Therefore, $\ker(T-\lambda_i\id)^{p_i}\cap\img(T-\lambda_i\id)^{p_i}=\{\0\}$.

    Since the two vector spaces are disjoint, they are also independent.
    Therefore,
    \begin{align*}
        &\dim\left(\ker(T-\lambda_i\id)^{p_i}\oplus\img(T-\lambda_i\id)^{p_i}\right)\\
        =&\dim\ker(T-\lambda_i\id)^{p_i} + \dim\img(T-\lambda_i\id)^{p_i}\\
        =&\dim\vecspace{V}\\
    \end{align*}
    by the dimension theorem.
    Hence, $\ker(T-\lambda_i\id)^{p_i}\oplus\img(T-\lambda_i\id)^{p_i}=\vecspace{V}$.
\end{proof}

\begin{theorem}
    $\vecspace{S}_0\oplus\vecspace{S}_1\oplus\dots\oplus\vecspace{S}_{k-1}=\vecspace{V}$
\end{theorem}
\begin{proof}
    Note that $\{\0\}\oplus\vecspace{S}=\vecspace{S}$. So for $k=0$, we will define $\vecspace{S}_0\oplus\dots\oplus\vecspace{S}_{k-1}=\{\0\}$.

    If $k=0$, since for any $n$-dimensional vector space $\vecspace{V}$ where $n > 0$, there exists eigenvalues for $T$, we have $n=0$.
    Thus, $\vecspace{V}=\{\0\}$ and the theorem holds.

    Suppose the theorem is true for some $k-1\in\N$. \\
    Consider the subspace $\vecspace{V}'=\img(T-\lambda_{k-1}\id)^{p_{k-1}}$.
    Note that from Lemma \ref{lemma:img_sum_ker}, we have $\vecspace{V} = \vecspace{S}_{k-1} \oplus \vecspace{V}'$.\\
    It is closed under $T$, as $\forall\vec{v}\in\vecspace{V}'$, there exist $\vec{u}\in\vecspace{V}$ s.t.
    \begin{equation*}
        \vec{v}=(T-\lambda_{k-1}\id)^{p_{k-1}}\vec{u}
    \end{equation*}
    Hence,
    \begin{equation*}
        T\vec{v} = T(T-\lambda_{k-1}\id)^{p_{k-1}}\vec{u} = (T-\lambda_{k-1}\id)^{p_{k-1}}(T\vec{u})\in\vecspace{V}'
    \end{equation*}

    Since $\vecspace{V}'\cap\vecspace{S}_{k-1}=\{\0\}$, $\spec_{\vecspace{V}'}T \subseteq \spec_{\vecspace{V}}T\setminus\{\lambda_{k-1}\}$.\\
    % $\forall i\in[k-1]$, let $\0\neq\vec{x}\in\ker(T-\lambda_i\id)\subseteq\vecspace{S}_i$.
    $\forall i\in[k-1]$, since $T-\lambda_{k-1}\id$ is a bijection on $\vecspace{S}_i$ (Lemma \ref{lemma:bijection}),
    \begin{equation*}
        \vecspace{S}_i=(T-\lambda_{k-1}\id)^{p_{k-1}}(\vecspace{S}_i)\subseteq\img(T-\lambda_{k-1}\id)^{p_{k-1}}=\vecspace{V}'
    \end{equation*}

    Hence, there are exactly $k-1$ eigen values of $T$ in $\vecspace{V}'$.
    We can apply induction hypothesis on $\vecspace{V}'$ to get
    \begin{equation*}
        \vecspace{V}' = \vecspace{S}_0\oplus\vecspace{S}_1\oplus\dots\oplus\vecspace{S}_{k-2}
    \end{equation*}

    Therefore,
    \begin{equation*}
        \vecspace{V} = \vecspace{V}' \oplus \vecspace{S}_{k-1} = \vecspace{S}_0\oplus\vecspace{S}_1\oplus\dots\oplus\vecspace{S}_{k-2} \oplus \vecspace{S}_{k-1}
    \end{equation*}

    %  let $\0\neq\vec{x}\in\ker(T-\lambda_i\id)\subseteq\vecspace{S}_i$.

    % need lemma to show kernel is closed under T
    % Note that $\0\neq(T-\lambda_{k-1}\id)^{p_{k-1}}\vec{x}\in\img(T-\lambda_{k-1}\id)^{p_{k-1}}$ and $(T-\lambda_i\id)(T-\lambda_{k-1}\id)^{p_{k-1}}\vec{x}=\0$.
    % Thus, $\lambda_i\in\spec_{\vecspace{V}'}T$.\\
    % Therefore, we can apply the induction hypothesis on $\vecspace{V}'$. We get
    % \begin{equation*}
    %     \vecspace{V}' = \vecspace{S}'_0\oplus\vecspace{S}'_1\oplus\dots\oplus\vecspace{S}'_{k-2}
    % \end{equation*}
    % where
    % \begin{equation*}
    %     \vecspace{S}'_i = \extker_{\vecspace{V}'}(T-\lambda_i\id)
    % \end{equation*}

\end{proof}

% \begin{lemma}
%     \label{lemma:construction_raw}
%     There exists linear independent subspaces $\vecspace{S_0}, \vecspace{S_1}, \dots, \vecspace{S_{k-1}}$ such that $\vecspace{V}=\vecspace{S_0} \oplus \vecspace{S_1} \oplus \dots \oplus \vecspace{S_{k-1}}$ and
%     $\vecspace{S_i} \subseteq \extker(T-\lambda_i \id)\sin$
% \end{lemma}
% \begin{proof}
%     We can construct such a list of subspaces recursively.
%     Let $\vecspace{V_0} = \vecspace{V}$. Let $S_i$
% \end{proof}



\end{document}

